% 图 18.4 —— 由 `checks/emit_hand.py` 自动拼出（probe 精测 + prims2 图元分解）
%%
% ⚠ 这是**稿子不是成品**：跑 `checks/checkhand.py 18.4` 看指标 + 看叠合图。
%%
%% ⚠ 待核：
%   · 本图 probe 报出阴影族 1 个 —— **本脚本不生成阴影**（要照原书手写）；prims2 的 strokes 里可能含阴影线，核对时留意别画重
%   · 可疑字形 1 项（空心圈 ⊙ 常被认成字母 o）—— 见文件末
%%
\documentclass[12pt]{standalone}
\usepackage{fontspec}
\setsansfont{Arial}
\newfontfamily\mfallback{DejaVu Sans}
\usepackage{tikz}
\usetikzlibrary{arrows.meta}
\begin{document}
\begin{tikzpicture}[x=1pt, y=-1pt, line width=4.0pt, line cap=round, line join=round]

  %% ════ 填充区（prims2 的 fills）════

  %% ════ 直线段（prims2 的 strokes）════
  \draw[line width=5.0pt] (900.0,213.0) -- (911.0,201.0);
  \draw[line width=4.0pt] (918.0,202.0) -- (911.0,243.1);
  \draw[line width=4.0pt, -{Triangle[width=10.8pt,length=16.2pt]}] (402.0,198.0) -- (496.0,196.0);
  \draw[line width=4.0pt] (937.0,212.0) -- (927.0,201.0);
  \draw[line width=6.0pt] (920.0,200.0) -- (926.0,200.0);
  \draw[line width=4.0pt] (231.0,197.0) -- (77.0,197.0);
  \draw[line width=9.0pt] (75.0,197.0) -- (25.0,197.0);
  \draw[line width=5.0pt] (912.0,201.0) -- (917.0,201.0);

  %% ════ 曲线（prims2 的 curves —— 三段式，坐标同为 (y,x)）════
  \draw[line width=5.0pt] (911.0,243.1) .. controls (892.6,326.5) and (795.0,386.3) .. (713.0,356.0);
  \draw[line width=5.0pt] (713.0,356.0) .. controls (689.7,351.6) and (669.6,337.6) .. (650.8,322.8);
  \draw[line width=5.0pt] (650.8,322.8) .. controls (605.1,281.4) and (588.3,219.5) .. (600.0,159.3);
  \draw[line width=5.0pt] (600.0,159.3) .. controls (617.6,93.1) and (678.6,34.2) .. (750.0,34.0);
  \draw[line width=8.0pt] (919.0,199.0) .. controls (917.6,156.8) and (904.0,116.6) .. (876.4,84.4);
  \draw[line width=8.0pt] (876.4,84.4) .. controls (843.2,50.5) and (799.7,33.8) .. (753.0,34.0);
  \draw[line width=5.0pt] (762.0,15.0) .. controls (755.5,18.7) and (750.6,24.9) .. (752.0,33.0);
  \draw[line width=4.0pt] (34.0,217.0) .. controls (21.2,219.6) and (19.0,206.2) .. (24.0,197.0);
  \draw[line width=5.0pt] (66.0,178.0) .. controls (73.4,179.8) and (78.0,189.4) .. (76.0,196.0);
  \draw[line width=4.0pt] (911.0,200.0) .. controls (903.8,197.7) and (897.0,197.5) .. (900.0,187.0);
  \draw[line width=4.0pt] (231.0,197.0) .. controls (233.4,189.5) and (228.1,180.6) .. (221.0,177.0);
  \draw[line width=5.0pt] (231.0,197.0) .. controls (233.3,203.7) and (227.0,211.7) .. (221.0,215.0);
  \draw[line width=6.0pt] (936.0,188.0) .. controls (939.9,195.1) and (934.9,199.8) .. (928.0,200.0);
  \draw[line width=5.0pt] (751.0,35.0) .. controls (750.3,42.6) and (753.3,51.6) .. (761.0,54.0);
  \draw[line width=4.5pt] (33.0,179.0) .. controls (21.4,174.8) and (19.1,188.7) .. (24.0,196.0);
  \draw[line width=5.0pt] (76.0,198.0) .. controls (77.4,205.8) and (72.2,212.2) .. (66.0,216.0);

  %% ════ 圆 / 椭圆（probe 精测）════
  \draw (757.8,199.1) circle (162.9);

  %% ════ 实心点 ════
  % （略）(480.5,192.5) r=4.5 落在箭头头上 —— 已由箭头头代替

  %% ════ 标签（probe 直接给的 \node 行）════
  %% ⚠ 全部补了 fill=white：文字在最上层，否则与曲线交叉干扰
     \node[anchor=center,inner sep=0pt,fill=white,font=\fontsize{27.1}{33.9}\selectfont] at (885.5,60.5) {\textsf{\textbf{\textit{f}}}};   % box(880,49,10,22)
     \node[anchor=center,inner sep=0pt,fill=white,font=\fontsize{30.7}{38.4}\selectfont] at (899.5,63.5) {\textsf{\textbf{(}}};   % box(895,49,8,28)
     \node[anchor=center,inner sep=0pt,fill=white,font=\fontsize{29.6}{37.0}\selectfont] at (917.0,61.2) {\textsf{\textbf{\textit{U}}}};   % box(907,49,20,22)
     \node[anchor=center,inner sep=0pt,fill=white,font=\fontsize{30.7}{38.4}\selectfont] at (935.5,63.5) {\textsf{\textbf{)}}};   % box(930,49,8,28)
     \node[anchor=center,inner sep=0pt,fill=white,font=\fontsize{29.6}{37.0}\selectfont] at (55.0,162.2) {\textsf{\textbf{\textit{U}}}};   % box(45,150,20,22)
     \node[anchor=center,inner sep=0pt,fill=white,font=\fontsize{29.1}{36.4}\selectfont] at (449.0,174.0) {\textsf{\textbf{\textit{f}}}};   % box(443,162,11,23)
     \node[anchor=center,inner sep=0pt,fill=white,font=\fontsize{29.8}{37.3}\selectfont] at (959.0,204.5) {\textsf{\textbf{\textit{P}}}};   % box(948,192,18,23)
     \node[anchor=center,inner sep=0pt,fill=white,font=\fontsize{31.1}{38.8}\selectfont] at (28.3,243.8) {\textsf{\textbf{6}}};   % box(20,232,15,22)
     \node[anchor=center,inner sep=0pt,fill=white,font=\fontsize{23.7}{29.7}\selectfont] at (73.8,240.4) {\textsf{\textbf{1}}};   % box(67,232,8,16)
     \node[anchor=center,inner sep=0pt,fill=white,font=\fontsize{30.4}{38.0}\selectfont] at (228.5,243.0) {\textsf{\textbf{1}}};   % box(220,232,10,21)
     \node[anchor=center,inner sep=0pt,fill=white,font=\fontsize{37.2}{46.4}\selectfont] at (80.5,257.6) {{\mfallback\textbf{−}}};   % box(63,249,17,6)
     \node[anchor=center,inner sep=0pt,fill=white,font=\fontsize{22.7}{28.3}\selectfont] at (72.1,265.6) {\textsf{\textbf{6}}};   % box(66,257,11,16)

  % 钉死画布：否则超出的图元/控制点会把页面撑大，1:1 回比就失准
  \pgfresetboundingbox
  \path[use as bounding box] (0,0) rectangle (975,374);
\end{tikzpicture}
\end{document}

% ⚠ 待核清单（probe 拿不准，人眼过一遍）：
%   · 未识别/跳过：⑤ 跳过/未识别的字形
